% Part: normal-modal-logic
% Chapter: axioms-systems
% Section: duals

\documentclass[../../../include/open-logic-section]{subfiles}

\begin{document}

\olfileid{nml}{prf}{dua}

\olsection{द्वयी \usetoken{P}{formula}}

\begin{defn}\ollabel{def:duals}
  \Ax{T}, \Ax{B}, \Ax{4} आणि \Ax{5} ही
  !!{formula}s आहेत; त्यांपैकी प्रत्येकाला
  \emph{द्वयी} सूत्र आहे.
  ते तळसूचक हिऱ्याने पुढीलप्रमाणे दर्शवले जाते:
  \begin{align}
    \tag{\Ax{T_\Diamond}} p & \lif \Diamond p\\
    \tag{\Ax{B_\Diamond}} \Diamond\Box p & \lif p\\
    \tag{\Ax{4_\Diamond}} \Diamond\Diamond p & \lif \Diamond p\\
    \tag{\Ax{5_\Diamond}} \Diamond\Box p & \lif \Box p
    \end{align}
\end{defn}

वरील द्वयी !!{formula}s पैकी प्रत्येक सूत्र संबंधित
!!{formula} मधील $p$ च्या जागी $\lnot p$ ठेवून,
प्रतिपरिवर्तन करून,
$\lnot\Box\lnot$ च्या जागी $\Diamond$ आणि
$\lnot\Diamond\lnot$ च्या जागी~$\Box$ ठेवून मिळते.
\Ax{D}, म्हणजे $\Box!A \lif \Diamond!A$, हे या अर्थाने
स्वतःचेच द्वयी सूत्र आहे.

\begin{prop}\ollabel{prop:dualsys}
  \olref[nml][prf][dua]{def:duals} मधील प्रत्येक
  !!{formula}~$!A$ साठी:
  $\Log{K}!A = \Log{K}!A_{\Diamond}$.
\end{prop}
\begin{proof}
  सराव.
\end{proof}
\begin{prob}
  \olref[nml][prf][dua]{prop:dualsys} सिद्ध करा.
\end{prob}

\end{document}
